\chapter{Trie}
\chapterepigraph{A trie trades memory for a kind of search no comparison-based
structure can offer: it does not ask ``is this string less than or
greater than that one'' -- it asks ``what character comes next,'' one
position at a time, branching only where words actually disagree.}

\section{What Problem Does a Trie Actually Solve?}
\label{sec:trie-intro}

A \textbf{trie} (prefix tree) stores a collection of strings (or, viewed
differently, a collection of fixed-length binary numbers) as paths from a
shared root, where every node represents one character (or one bit)
consumed so far, and nodes are shared across all strings that agree on a
prefix. Two strings that share a $k$-character prefix share the first $k$
nodes of their paths through the trie; they only diverge into separate
branches at the first character where they differ.

\begin{ideabox}[Why This Beats a Hash Set or a BST for Prefix Questions]
A hash set answers ``is this exact string present'' in $O(\text{length})$,
but cannot answer ``how many stored strings start with this prefix''
without scanning every stored string. A BST of strings can answer prefix
questions via range queries, but each comparison still costs
$O(\text{length})$ in the worst case, and the tree's shape depends on
insertion order. A trie answers \textbf{both} exact-match and
prefix-match queries in $O(\text{length of the query})$, completely
independent of how many strings are stored or in what order they were
inserted -- because shared prefixes are physically shared as the same
path, prefix counting is just a matter of annotating nodes along that
shared path with counts as strings are inserted.
\end{ideabox}

\subsection*{How to Recognise a Trie Problem from Its Wording}

\begin{patternbox}[Recurring Trie Keywords]
\begin{itemize}
  \item \textbf{``Insert/search words,'' ``starts with,'' ``prefix''} --
        a direct, textbook trie: one node per character, a marker at
        nodes where a complete word ends.
  \item \textbf{``Count words equal to / starting with a given
        string,'' or ``erase a word''} -- the same trie, with each node
        additionally annotated with a small counter (how many inserted
        words pass through, or end at, this node), updated incrementally
        on insert/erase.
  \item \textbf{``Longest prefix that is itself always a complete,
        previously-seen word'' or ``every prefix of this word must also
        be a complete word''} -- a trie where the search/build step
        explicitly checks, at every node along a path, whether that node
        marks the end of \emph{some} inserted word, not just whether the
        final node does.
  \item \textbf{``Maximum XOR'' between numbers, or ``maximum XOR with
        an element from an array,'' for a query''} -- a \textbf{binary
        trie}: treat each number as a fixed-length sequence of bits
        (e.g.\ $32$ bits), insert each number bit by bit from the most
        significant bit downward, and answer a maximum-XOR query by
        greedily choosing, at each bit position, the branch that
        \emph{disagrees} with the query number's corresponding bit
        (since XOR-ing differing bits produces a $1$, maximizing the
        result).
  \item \textbf{``Number of distinct substrings''} -- build a trie of
        \emph{every suffix} of the string; the number of distinct
        substrings equals the number of nodes created during this
        construction (every trie node corresponds to exactly one
        distinct substring, the one spelled out by the path from the
        root to that node).
\end{itemize}
\end{patternbox}

Each section in this chapter follows the established structure: problem
statement and keywords; the reasoning that identifies the trie variant
needed and why it is correct; the algorithm in words and pseudocode; a
hand-traced dry run; complete compilable C++ code; and a time/space
complexity analysis.

\newpage

\section{Implement Trie (Insert, Search, StartsWith)}
\label{sec:trie-implement}

\problemstatement{Implement a trie supporting $\textit{insert}(word)$,
$\textit{search}(word)$ (returns true only if $word$ was previously
inserted as a \emph{complete} word), and $\textit{startsWith}(prefix)$
(returns true if any inserted word begins with $prefix$).}

\begin{patternbox}
This is the foundational trie, with no special twist -- the rest of this
chapter builds directly on the node structure and walk logic derived
here. The one subtlety worth flagging immediately: $\textit{search}$ and
$\textit{startsWith}$ use the \emph{same} walk, differing only in what
they check once the walk finishes.
\end{patternbox}

\subsection*{Understanding the problem carefully}

Each trie node represents one character position along some prefix, and
holds up to $26$ child pointers (one per lowercase letter, say), plus a
boolean flag marking whether a complete inserted word \emph{ends} exactly
at this node. Walking a string through the trie, one character at a
time, either follows existing child pointers (if a previous word shares
this prefix) or creates new nodes (during insertion, for the first word
to diverge here).

\begin{ideabox}[One Walk, Three Different Endings]
$\textit{insert}(word)$: starting at the root, for each character,
create the corresponding child if it does not yet exist, then move into
it; after the last character, mark the final node's end-of-word flag.
$\textit{search}(word)$: walk the same way, but if any required child is
missing, return false immediately; after the last character, return
whether the final node's end-of-word flag is set (not just whether the
path exists -- $word$ might only be a \emph{prefix} of something
inserted, not itself a complete inserted word). $\textit{startsWith}
(prefix)$: identical walk, but return true as soon as the path exists in
full, regardless of the end-of-word flag.
\end{ideabox}

\subsection*{Why search must check the end-of-word flag but startsWith must not}

This is the one place a careless implementation typically goes wrong: if
only ``apple'' has been inserted, then $\textit{search}(\text{"app"})$
must return \texttt{false} (``app'' was never inserted as a complete
word, even though it is a valid path in the trie), while
$\textit{startsWith}(\text{"app"})$ must return \texttt{true} (something
inserted does begin with ``app''). The trie's path existing is necessary
but not sufficient for $\textit{search}$ -- it additionally needs the
explicit end-of-word marker, which is exactly the piece of information
$\textit{startsWith}$ does not need to check at all.

\subsection*{Algorithm}

\begin{enumerate}
  \item \textsc{Insert}($word$): $\textit{node} \gets \textit{root}$.
        For each character $c$ in $word$: if $\textit{node}$ has no
        child $c$, create one; move $\textit{node}$ into child $c$.
        Mark $\textit{node}.\textit{isEnd} \gets \text{true}$.
  \item \textsc{Search}($word$): walk as above; if any required child
        is missing, return false; otherwise return
        $\textit{node}.\textit{isEnd}$.
  \item \textsc{StartsWith}($prefix$): walk as above; return true if the
        full walk succeeds (regardless of $\textit{isEnd}$), false if any
        child is missing.
\end{enumerate}

\subsection*{Dry run}

\dryrun{Insert \texttt{"apple"}. Query \textit{search}(\texttt{"app"}),
\textit{startsWith}(\texttt{"app"}), \textit{search}(\texttt{"apple"}).}

\begin{itemize}
  \item Insert \texttt{"apple"}: creates a path of $5$ nodes,
        $a\to p\to p\to l\to e$, marking only the final ($e$) node's
        $\textit{isEnd}$ flag.
  \item $\textit{search}(\texttt{"app"})$: path $a\to p\to p$ exists,
        but that node's $\textit{isEnd}$ is false (only the $e$ node at
        the end of ``apple'' has it set). Returns \texttt{false}.
  \item $\textit{startsWith}(\texttt{"app"})$: path exists, so returns
        \texttt{true}, regardless of the flag.
  \item $\textit{search}(\texttt{"apple"})$: path exists and the final
        node's $\textit{isEnd}$ is true. Returns \texttt{true}.
\end{itemize}

\subsection*{C++ implementation}

\begin{lstlisting}
#include <bits/stdc++.h>
using namespace std;

struct TrieNode {
    TrieNode* children[26] = {nullptr};
    bool isEnd = false;
};

class Trie {
    TrieNode* root;

public:
    Trie() { root = new TrieNode(); }

    void insert(string word) {
        TrieNode* node = root;
        for (char c : word) {
            int idx = c - 'a';
            if (!node->children[idx]) node->children[idx] = new TrieNode();
            node = node->children[idx];
        }
        node->isEnd = true;
    }

    bool search(string word) {
        TrieNode* node = root;
        for (char c : word) {
            int idx = c - 'a';
            if (!node->children[idx]) return false;
            node = node->children[idx];
        }
        return node->isEnd;
    }

    bool startsWith(string prefix) {
        TrieNode* node = root;
        for (char c : prefix) {
            int idx = c - 'a';
            if (!node->children[idx]) return false;
            node = node->children[idx];
        }
        return true;
    }
};

int main() {
    Trie trie;
    trie.insert("apple");
    cout << trie.search("app") << endl;       // 0 (false)
    cout << trie.startsWith("app") << endl;   // 1 (true)
    cout << trie.search("apple") << endl;     // 1 (true)
    return 0;
}
\end{lstlisting}

\begin{complexitybox}
\textbf{Time Complexity.} $O(L)$ for each operation, where $L$ is the
length of the word/prefix involved -- independent of how many other
words are stored in the trie.
\textbf{Space Complexity.} $O(\Sigma \cdot N)$ in the worst case, where
$N$ is the total number of characters inserted across all words and
$\Sigma$ is the alphabet size (each node reserves space for all $26$
possible children, whether used or not); shared prefixes reduce this in
practice.
\end{complexitybox}

\begin{takeawaybox}
The entire chapter's machinery is this one node structure (a fixed-size
child array plus an end marker) and this one walk (follow or create
children, one character at a time). Every later section either adds more
annotations to each node (counts, in
Section~\ref{sec:trie-implement-2}) or changes what ``character'' means
(bits, in Sections~\ref{sec:trie-max-xor}--\ref{sec:trie-max-xor-query}).
\end{takeawaybox}

\section{Implement Trie II (Count and Erase)}
\label{sec:trie-implement-2}

\problemstatement{Implement a trie supporting $\textit{insert}(word)$,
$\textit{countWordsEqualTo}(word)$ (how many times $word$ has been
inserted, counting duplicates), $\textit{countWordsStartingWith}(prefix)$
(how many inserted words, counting duplicates, begin with $prefix$), and
$\textit{erase}(word)$ (remove one occurrence of a previously inserted
$word$).}

\begin{patternbox}
The keyword shift from Section~\ref{sec:trie-implement} is from a
yes/no question (\emph{is this word present}) to a \textbf{counting}
question (\emph{how many times}, including duplicates and including
prefix counts), plus \textbf{removal}. The fix is to annotate every node,
not just mark a single boolean: track how many inserted words
\textbf{pass through} this node (as part of some longer word, or ending
exactly here), and how many \textbf{end exactly} at this node.
\end{patternbox}

\subsection*{Understanding the problem carefully}

Every node along the path of an inserted word is, by definition, a
prefix of that word -- so if we increment a $\textit{countPrefix}$
counter at \emph{every} node visited during an insertion, that counter at
any given node ends up holding exactly the number of inserted words
(counting duplicates) that share this far as a common prefix. Separately,
incrementing an $\textit{countEnd}$ counter only at the \emph{final} node
of each inserted word tracks exactly how many times that complete word
itself was inserted.

\begin{ideabox}[Two Counters per Node: Prefix-Passes and Exact-Ends]
$\textit{insert}(word)$: walk the word's path, creating nodes as needed
(as in Section~\ref{sec:trie-implement}), and increment
$\textit{countPrefix}$ at \emph{every} node visited along the way
(including the final one); increment $\textit{countEnd}$ only at the
final node. $\textit{countWordsEqualTo}(word)$: walk the path; if it
exists in full, return the final node's $\textit{countEnd}$ (else $0$).
$\textit{countWordsStartingWith}(prefix)$: walk the path; if it exists in
full, return the final node's $\textit{countPrefix}$ (else $0$).
$\textit{erase}(word)$: walk the path (assuming $word$ was indeed
previously inserted), \textbf{decrementing} $\textit{countPrefix}$ at
every node visited, and decrementing $\textit{countEnd}$ at the final
node.
\end{ideabox}

\subsection*{Why two separate counters are needed, not just one}

A single counter cannot distinguish ``$\textit{prefix}$ is itself a
complete inserted word $k$ times'' from ``$\textit{prefix}$ is merely a
prefix of $k$ longer inserted words.'' For example, after inserting
``app'' once and ``apple'' once, the node at the end of ``app'' has
$\textit{countPrefix}=2$ (both words pass through it) but
$\textit{countEnd}=1$ (only ``app'' itself ends there) --
$\textit{countWordsEqualTo}(\texttt{"app"})$ must report $1$, while
$\textit{countWordsStartingWith}(\texttt{"app"})$ must report $2$, and
no single counter can serve both queries correctly.

\subsection*{Algorithm}

\begin{enumerate}
  \item \textsc{Insert}($word$): walk/create the path; increment
        $\textit{countPrefix}$ at every visited node; increment
        $\textit{countEnd}$ at the final node.
  \item \textsc{CountWordsEqualTo}($word$): walk the path; return the
        final node's $\textit{countEnd}$, or $0$ if the path does not
        fully exist.
  \item \textsc{CountWordsStartingWith}($prefix$): walk the path; return
        the final node's $\textit{countPrefix}$, or $0$ if the path does
        not fully exist.
  \item \textsc{Erase}($word$): walk the path; decrement
        $\textit{countPrefix}$ at every visited node; decrement
        $\textit{countEnd}$ at the final node.
\end{enumerate}

\subsection*{Dry run}

\dryrun{Insert \texttt{"apple"} twice, insert \texttt{"app"} once. Query
$\textit{countWordsEqualTo}(\texttt{"apple"})$,
$\textit{countWordsStartingWith}(\texttt{"app"})$. Then erase
\texttt{"apple"} once and repeat both queries.}

\begin{itemize}
  \item After both insertions: node at \texttt{"app"} has
        $\textit{countPrefix}=3$ (all three insertions pass through
        it), $\textit{countEnd}=1$ (only the single ``app'' insertion
        ends there). Node at \texttt{"apple"} has
        $\textit{countPrefix}=2$, $\textit{countEnd}=2$.
  \item $\textit{countWordsEqualTo}(\texttt{"apple"}) = 2$.
        $\textit{countWordsStartingWith}(\texttt{"app"}) = 3$.
  \item Erase \texttt{"apple"} once: decrement $\textit{countPrefix}$
        along the whole path (node at \texttt{"app"}: $3\to2$; node at
        \texttt{"apple"}: $2\to1$), and $\textit{countEnd}$ at the final
        node ($2\to1$).
  \item Now $\textit{countWordsEqualTo}(\texttt{"apple"}) = 1$.
        $\textit{countWordsStartingWith}(\texttt{"app"}) = 2$.
\end{itemize}

\subsection*{C++ implementation}

\begin{lstlisting}
#include <bits/stdc++.h>
using namespace std;

struct TrieNode {
    TrieNode* children[26] = {nullptr};
    int countPrefix = 0;
    int countEnd = 0;
};

class TrieII {
    TrieNode* root;

public:
    TrieII() { root = new TrieNode(); }

    void insert(string word) {
        TrieNode* node = root;
        for (char c : word) {
            int idx = c - 'a';
            if (!node->children[idx]) node->children[idx] = new TrieNode();
            node = node->children[idx];
            node->countPrefix++;
        }
        node->countEnd++;
    }

    int countWordsEqualTo(string word) {
        TrieNode* node = root;
        for (char c : word) {
            int idx = c - 'a';
            if (!node->children[idx]) return 0;
            node = node->children[idx];
        }
        return node->countEnd;
    }

    int countWordsStartingWith(string prefix) {
        TrieNode* node = root;
        for (char c : prefix) {
            int idx = c - 'a';
            if (!node->children[idx]) return 0;
            node = node->children[idx];
        }
        return node->countPrefix;
    }

    void erase(string word) {
        TrieNode* node = root;
        for (char c : word) {
            int idx = c - 'a';
            if (!node->children[idx]) return; // not present
            node = node->children[idx];
            node->countPrefix--;
        }
        node->countEnd--;
    }
};

int main() {
    TrieII trie;
    trie.insert("apple");
    trie.insert("apple");
    trie.insert("app");

    cout << trie.countWordsEqualTo("apple") << endl;        // 2
    cout << trie.countWordsStartingWith("app") << endl;      // 3

    trie.erase("apple");
    cout << trie.countWordsEqualTo("apple") << endl;        // 1
    cout << trie.countWordsStartingWith("app") << endl;      // 2
    return 0;
}
\end{lstlisting}

\begin{complexitybox}
\textbf{Time Complexity.} $O(L)$ for every operation.
\textbf{Space Complexity.} $O(\Sigma \cdot N)$, as in
Section~\ref{sec:trie-implement}, with two integer counters added per
node.
\end{complexitybox}

\begin{takeawaybox}
Whenever a trie problem needs counts rather than a single boolean,
annotate every node with the counters the queries actually need --
typically ``how many words pass through here'' and/or ``how many words
end exactly here'' -- and update them incrementally on insert and erase.
This is a direct generalization of the single $\textit{isEnd}$ flag from
Section~\ref{sec:trie-implement}.
\end{takeawaybox}

\section{Longest String with All Prefixes}
\label{sec:trie-longest-word-prefixes}

\problemstatement{Given an array of words, find the \textbf{longest}
word such that \textbf{every prefix} of that word (of every length from
$1$ up to its full length) is also present, as a complete word, somewhere
in the array. If there are ties in length, return the lexicographically
smallest such word. If no word qualifies (not even single-character
words), return the empty string.}

\begin{patternbox}
\emph{``Every prefix of this word must also be a complete word in the
collection''} is the keyword that distinguishes this from
Section~\ref{sec:trie-implement}'s plain $\textit{startsWith}$ check: we
are not asking whether a path exists, but whether \textbf{every node
along that entire path} happens to be an end-of-word marker, not just the
final node.
\end{patternbox}

\subsection*{Understanding the problem carefully}

Insert every word into a trie, exactly as in
Section~\ref{sec:trie-implement}, marking end-of-word at each word's
final node. Then, for each word, walk its path through the trie and check
\emph{every} node visited (not just the last one) for the end-of-word
marker. A word qualifies if and only if this check passes at every
single position along its path -- meaning the 1-character prefix, the
2-character prefix, and so on, are all themselves complete words that
were inserted into the trie.

\begin{ideabox}[Check Every Node Along the Path, Not Just the Last]
For each candidate word $w$, walk the trie character by character. At
\emph{every} node visited (after consuming each character), check that
node's end-of-word flag. If any node along the way fails this check, $w$
does not qualify -- stop checking this word immediately. If the entire
walk completes with every node passing the check, $w$ qualifies; compare
it against the best answer found so far (preferring longer length, and
lexicographically smaller on ties).
\end{ideabox}

\subsection*{Why checking the whole path (not just the destination) is the correct test}

The problem's requirement is explicitly about \emph{every prefix}, not
just the word itself -- a word might be present in the array (so its
final node is marked end-of-word) while one of its shorter prefixes was
\emph{never} separately inserted as its own word, which would still
disqualify it. Since every node on a word's path through the trie
corresponds to exactly one specific prefix of that word, checking the
end-of-word flag at every node along the path is precisely equivalent to
checking, for every prefix length, whether that prefix was itself
inserted as a complete word.

\subsection*{Algorithm}

\begin{enumerate}
  \item Insert every word into a trie (as in
        Section~\ref{sec:trie-implement}).
  \item For each word $w$ (processed, say, in sorted order so that ties
        are resolved automatically by encountering the lexicographically
        smaller one first):
  \begin{enumerate}
    \item Walk $w$'s path; at each node, check the end-of-word flag; if
          any check fails, discard $w$.
    \item If $w$ survives the full walk and is longer than the current
          best answer (or ties in length but is lexicographically
          smaller, naturally handled by sorted-order processing), update
          the best answer.
  \end{enumerate}
  \item Return the best answer found (or empty string if none
        qualified).
\end{enumerate}

\subsection*{Dry run}

\dryrun{Take $\textit{words} = [\texttt{"a"}, \texttt{"banana"},
\texttt{"app"}, \texttt{"appl"}, \texttt{"ap"}, \texttt{"apply"},
\texttt{"apple"}]$.}

Sorted: \texttt{"a", "ap", "app", "appl", "apple", "apply", "banana"}.

\begin{itemize}
  \item \texttt{"a"}: path is just the node for \texttt{a}; it is marked
        end-of-word (since \texttt{"a"} itself was inserted). Qualifies,
        length $1$. Best so far: \texttt{"a"}.
  \item \texttt{"ap"}: nodes for \texttt{a}, \texttt{ap}; both marked
        (since \texttt{"a"} and \texttt{"ap"} were both inserted).
        Qualifies, length $2$. Best so far: \texttt{"ap"}.
  \item \texttt{"app"}: nodes for \texttt{a}, \texttt{ap}, \texttt{app}
        all marked. Qualifies, length $3$. Best so far: \texttt{"app"}.
  \item \texttt{"appl"}: all four prefix nodes marked. Qualifies, length
        $4$. Best so far: \texttt{"appl"}.
  \item \texttt{"apple"}: all five prefix nodes marked (every prefix
        from \texttt{"a"} to \texttt{"apple"} was inserted). Qualifies,
        length $5$. Best so far: \texttt{"apple"}.
  \item \texttt{"apply"}: prefix \texttt{"appl"} is marked, but at the
        final node (\texttt{"apply"}), it is marked too -- but wait, we
        need \emph{every} prefix; \texttt{"apply"} shares
        \texttt{"a","ap","app","appl"} with \texttt{"apple"} (all
        marked), and \texttt{"apply"} itself was inserted, so it also
        qualifies at length $5$ -- tied with \texttt{"apple"}. Since
        \texttt{"apple"} was already recorded and is lexicographically
        smaller than \texttt{"apply"}, we keep \texttt{"apple"}.
  \item \texttt{"banana"}: the node for \texttt{b} is not marked
        end-of-word as a 1-letter prefix unless \texttt{"b"} itself was
        inserted -- it was not. Disqualified immediately.
\end{itemize}

Final answer: \texttt{"apple"}.

\subsection*{C++ implementation}

\begin{lstlisting}
#include <bits/stdc++.h>
using namespace std;

struct TrieNode {
    TrieNode* children[26] = {nullptr};
    bool isEnd = false;
};

class Trie {
public:
    TrieNode* root = new TrieNode();

    void insert(string& word) {
        TrieNode* node = root;
        for (char c : word) {
            int idx = c - 'a';
            if (!node->children[idx]) node->children[idx] = new TrieNode();
            node = node->children[idx];
        }
        node->isEnd = true;
    }

    bool allPrefixesExist(string& word) {
        TrieNode* node = root;
        for (char c : word) {
            int idx = c - 'a';
            if (!node->children[idx] || !node->children[idx]->isEnd) {
                return false;
            }
            node = node->children[idx];
        }
        return true;
    }
};

string longestCompleteString(vector<string>& words) {
    Trie trie;
    for (string& w : words) trie.insert(w);

    sort(words.begin(), words.end());

    string best = "";
    for (string& w : words) {
        if (trie.allPrefixesExist(w) && w.size() > best.size()) {
            best = w;
        }
    }
    return best;
}

int main() {
    vector<string> words = {"a","banana","app","appl","ap","apply","apple"};
    cout << longestCompleteString(words) << endl;
    return 0;
}
\end{lstlisting}

\begin{complexitybox}
\textbf{Time Complexity.} Inserting all words costs $O(N)$, where $N$ is
the total number of characters across all words. Sorting costs
$O(N \log n)$ (comparing $n$ words). Checking every word's path costs
another $O(N)$ total. Overall:
\[
  O(N \log n).
\]
\textbf{Space Complexity.} $O(\Sigma \cdot N)$ for the trie.
\end{complexitybox}

\begin{takeawaybox}
Whenever a problem's condition is phrased as ``true for every prefix,''
the trie check needs to inspect every node along a path, not just the
final destination -- a small but easy-to-miss generalization of plain
$\textit{search}$ from Section~\ref{sec:trie-implement}.
\end{takeawaybox}

\section{Number of Distinct Substrings in a String}
\label{sec:trie-distinct-substrings}

\problemstatement{Given a string $s$ of length $n$, count the number of
\textbf{distinct} (non-empty) substrings of $s$.}

\begin{patternbox}
\emph{``Count distinct substrings''} reframes naturally once you notice
that every substring of $s$ is, by definition, a \textbf{prefix of some
suffix} of $s$. Inserting \emph{every suffix} of $s$ into a trie
therefore inserts every substring of $s$ as some path from the root --
and because a trie physically \textbf{shares} identical prefixes as the
same nodes, the number of \textbf{distinct} substrings is exactly the
number of \textbf{distinct nodes} the trie ends up containing.
\end{patternbox}

\subsection*{Understanding the problem carefully}

Consider inserting all $n$ suffixes of $s$ (the suffix starting at index
$0$, the suffix starting at index $1$, \dots, the suffix starting at
index $n-1$) into a trie, one character at a time, exactly as in
Section~\ref{sec:trie-implement}. Every prefix of every suffix is some
substring of $s$, and every substring of $s$ is the prefix of at least
one suffix (specifically, the suffix that starts where the substring
starts) -- so this process inserts every substring of $s$, with
duplicates naturally \emph{not} creating new nodes, since inserting a
character that already has a corresponding child simply reuses that
existing node.

\begin{ideabox}[Count New Nodes Created, Not Total Insertions]
For each suffix of $s$ (there are $n$ of them, of lengths $n, n-1, \dots,
1$), walk it through the trie character by character. Whenever a
required child does not yet exist, create it \emph{and increment a
running counter of distinct substrings} (since a brand-new node means
this exact substring -- the path from root to this new node -- has never
been encountered before). When a required child already exists, simply
move into it without incrementing the counter (this exact substring was
already counted during an earlier suffix's insertion).
\end{ideabox}

\subsection*{Why counting newly created nodes exactly counts distinct substrings}

Each trie node corresponds to exactly one specific substring -- the
concatenation of characters along the path from the root to that node --
and two different nodes can never correspond to the same substring (a
trie's structure forbids two different paths spelling out the same
string, since shared prefixes are always merged into the same nodes).
Conversely, every substring of $s$ does correspond to \emph{some} node,
because it is a prefix of at least one inserted suffix. So there is a
perfect one-to-one correspondence between distinct substrings of $s$ and
nodes of the trie (excluding the root, which represents the empty
string) -- counting nodes created during insertion is therefore counting
distinct substrings directly, with no separate deduplication step needed.

\subsection*{Algorithm}

\begin{enumerate}
  \item Initialize an empty trie and $\textit{count} \gets 0$.
  \item For each starting index $i$ from $0$ to $n-1$ (each suffix
        $s[i..n-1]$):
  \begin{enumerate}
    \item Walk the suffix's characters from the trie's root. For each
          character: if the required child does not exist, create it
          and increment $\textit{count}$; move into the child.
  \end{enumerate}
  \item Return \textit{count}.
\end{enumerate}

\subsection*{Dry run}

\dryrun{Take $s = \texttt{"ab"}$. Suffixes: \texttt{"ab"},
\texttt{"b"}.}

\begin{itemize}
  \item Insert \texttt{"ab"}: \texttt{a} does not exist at the root --
        create it, $\textit{count}=1$. Move into \texttt{a}.
        \texttt{b} does not exist there -- create it, $\textit{count}=2$.
  \item Insert \texttt{"b"}: \texttt{b} does not exist at the root
        (this is a different node than the \texttt{b} created above,
        which was a child of \texttt{a}, not of the root) -- create it,
        $\textit{count}=3$.
\end{itemize}

Total distinct substrings: $3$ -- namely \texttt{"a"}, \texttt{"ab"},
\texttt{"b"}, matching the direct enumeration of all substrings of
\texttt{"ab"}.

\subsection*{C++ implementation}

\begin{lstlisting}
#include <bits/stdc++.h>
using namespace std;

struct TrieNode {
    TrieNode* children[26] = {nullptr};
};

int countDistinctSubstrings(string s) {
    TrieNode* root = new TrieNode();
    int n = s.size();
    int count = 0;

    for (int i = 0; i < n; i++) {
        TrieNode* node = root;
        for (int j = i; j < n; j++) {
            int idx = s[j] - 'a';
            if (!node->children[idx]) {
                node->children[idx] = new TrieNode();
                count++;
            }
            node = node->children[idx];
        }
    }
    return count;
}

int main() {
    cout << countDistinctSubstrings("ab") << endl;   // 3
    cout << countDistinctSubstrings("aaa") << endl;  // 3 (a, aa, aaa)
    return 0;
}
\end{lstlisting}

\begin{complexitybox}
\textbf{Time Complexity.} Inserting all $n$ suffixes, of total length
$n + (n-1) + \cdots + 1$, gives
\[
  O(n^2)
\]
character insertions in the worst case, each $O(1)$ amortized trie work.
(A suffix automaton or suffix array based approach can solve the same
problem in $O(n \log n)$ or $O(n)$, but the trie-of-all-suffixes
approach shown here is the direct, easy-to-derive $O(n^2)$ solution,
appropriate for moderate string lengths.)
\textbf{Space Complexity.} $O(n^2)$ in the worst case (one trie node per
distinct substring, and a string can have up to $O(n^2)$ distinct
substrings).
\end{complexitybox}

\begin{takeawaybox}
``Count distinct substrings'' is solved by recognizing that a trie's
node-sharing behavior \emph{is} a deduplication mechanism: insert
everything that could possibly create a duplicate (here, every suffix,
since every substring is a prefix of some suffix), and count only the
genuinely new nodes created, never the total insertion attempts.
\end{takeawaybox}

\section{Maximum XOR of Two Numbers in an Array}
\label{sec:trie-max-xor}

\problemstatement{Given an array $\textit{nums}$ of non-negative
integers, find the maximum value of $\textit{nums}[i] \oplus
\textit{nums}[j]$ over all pairs $i \ne j$ (where $\oplus$ denotes
bitwise XOR).}

\begin{patternbox}
\emph{``Maximum XOR''} between pairs of numbers is the keyword for a
\textbf{binary trie}: treat each number as a fixed-length sequence of
bits (most significant bit first) and insert it into a trie exactly as
Section~\ref{sec:trie-implement} inserts characters of a string -- except
the ``alphabet'' here has only two symbols, $0$ and $1$. Once every
number is in the trie, maximizing the XOR with a given number becomes a
greedy walk: at each bit position, XOR is maximized by making that bit
\textbf{differ}, so we always prefer the branch \emph{opposite} to the
current number's bit, whenever it exists.
\end{patternbox}

\subsection*{Understanding the problem carefully}

XOR-ing two bits produces $1$ exactly when they differ. To make the
overall XOR as large as possible, we want to maximize the \emph{most
significant} differing bit first (since higher bit positions dominate
the numeric value), then the next, and so on -- exactly the same
most-significant-first greedy logic used throughout binary search's
``binary search on the answer'' problems, but here applied bit by bit
through a trie rather than narrowing a numeric range.

\begin{ideabox}[Insert All Numbers as Bit-Paths, Then Greedily Walk for the Opposite Bit]
Insert every number into a binary trie, one bit at a time from the most
significant bit downward (using a fixed bit width, e.g.\ $32$, so all
numbers have aligned representations). For each number $x$ in the array,
walk the trie trying, at each bit position, to move into the child
representing the \emph{opposite} of $x$'s bit at that position; if that
child does not exist, settle for the only available child (the same
bit). Accumulate the resulting XOR value bit by bit as you walk, and
track the maximum such value found over all numbers $x$.
\end{ideabox}

\subsection*{Why the greedy opposite-bit choice is always optimal}

Because bit positions contribute to the final value with strictly
decreasing weight (the most significant bit contributes more than the
sum of all lower bits combined, for a fixed bit width), it is always
better to secure a $1$ at a higher bit position than to optimize any
combination of lower bits -- so greedily taking the opposite-bit branch
whenever it is available, at every level from the most significant bit
down, is guaranteed to produce the maximum possible XOR with the
existing numbers in the trie, by the same most-significant-digit-first
reasoning used whenever comparing fixed-width binary numbers.

\subsection*{Algorithm}

\begin{enumerate}
  \item Insert every number in $\textit{nums}$ into a binary trie, bit
        by bit from the most significant bit (say bit $31$) to the
        least significant (bit $0$).
  \item For each number $x$ in $\textit{nums}$:
  \begin{enumerate}
    \item Walk the trie from the root. At each bit position $b$ (from
          most to least significant): let $\textit{desired}$ be the
          opposite of $x$'s bit $b$. If the trie node has a child for
          $\textit{desired}$, move into it and set this bit of the
          running XOR result to $1$; otherwise, move into the only
          available child and set this bit to $0$.
    \item Update the global maximum with this number's resulting XOR
          value.
  \end{enumerate}
  \item Return the global maximum.
\end{enumerate}

\subsection*{Dry run}

\dryrun{Take $\textit{nums} = [3,10,5,25,2,8]$, using $5$-bit
representations for clarity ($3{=}00011$, $10{=}01010$, $5{=}00101$,
$25{=}11001$, $2{=}00010$, $8{=}01000$). Query the best XOR partner for
$5{=}00101$.}

\begin{itemize}
  \item Bit $4$ (value $16$) of $5$ is $0$; desired opposite is $1$. Only
        $25$ has bit $4=1$. Move there; XOR bit $4=1$.
  \item Bit $3$ (value $8$) of $5$ is $0$; desired $1$. Within this
        branch, only $25$ remains, and its bit $3=1$ -- matches desired.
        Move there; XOR bit $3=1$.
  \item Bit $2$ (value $4$) of $5$ is $1$; desired $0$. $25$'s bit
        $2=0$ -- matches. Move there; XOR bit $2=1$.
  \item Bit $1$ (value $2$) of $5$ is $0$; desired $1$. Only $25$'s
        branch remains here, and its bit $1=0$ -- not the desired
        branch, but it's the only one available. Move there; XOR bit
        $1=0$.
  \item Bit $0$ (value $1$) of $5$ is $1$; desired $0$. Only $25$'s
        bit $0=1$ available. Move there; XOR bit $0=0$.
\end{itemize}

Resulting XOR bits: $11100_2 = 28$. Indeed $5 \oplus 25 = 28$, the
maximum pairwise XOR for this array.

\subsection*{C++ implementation}

\begin{lstlisting}
#include <bits/stdc++.h>
using namespace std;

struct TrieNode {
    TrieNode* children[2] = {nullptr, nullptr};
};

class BinaryTrie {
    TrieNode* root = new TrieNode();
    static const int BITS = 32;

public:
    void insert(int num) {
        TrieNode* node = root;
        for (int b = BITS - 1; b >= 0; b--) {
            int bit = (num >> b) & 1;
            if (!node->children[bit]) node->children[bit] = new TrieNode();
            node = node->children[bit];
        }
    }

    int maxXorWith(int num) {
        TrieNode* node = root;
        int result = 0;
        for (int b = BITS - 1; b >= 0; b--) {
            int bit = (num >> b) & 1;
            int desired = 1 - bit;
            if (node->children[desired]) {
                result |= (1 << b);
                node = node->children[desired];
            } else {
                node = node->children[bit];
            }
        }
        return result;
    }
};

int findMaximumXOR(vector<int>& nums) {
    BinaryTrie trie;
    for (int x : nums) trie.insert(x);

    int maxXor = 0;
    for (int x : nums) {
        maxXor = max(maxXor, trie.maxXorWith(x));
    }
    return maxXor;
}

int main() {
    vector<int> nums = {3, 10, 5, 25, 2, 8};
    cout << "Maximum XOR: " << findMaximumXOR(nums) << endl;
    return 0;
}
\end{lstlisting}

\begin{complexitybox}
\textbf{Time Complexity.} Inserting all $n$ numbers costs $O(n \cdot
\textit{BITS})$. Querying the best partner for each of the $n$ numbers
costs another $O(n \cdot \textit{BITS})$. Overall:
\[
  O(n \cdot \textit{BITS}) = O(n),
\]
treating the fixed bit width ($32$) as a constant -- a dramatic
improvement over the $O(n^2)$ brute-force check of every pair.
\textbf{Space Complexity.} $O(n \cdot \textit{BITS})$ for the trie in the
worst case.
\end{complexitybox}

\begin{takeawaybox}
A binary trie is a character trie with a two-letter alphabet -- every
technique from Sections~\ref{sec:trie-implement}--\ref{sec:trie-distinct-substrings}
carries over unchanged, just with bits instead of letters. Whenever a
problem involves maximizing or counting based on XOR (or, more generally,
bitwise structure) across many numbers, a binary trie processed
most-significant-bit-first is almost always the intended technique.
\end{takeawaybox}

\section{Maximum XOR With an Element from Array (Queries)}
\label{sec:trie-max-xor-query}

\problemstatement{Given an array $\textit{nums}$ and a list of queries,
where each query is a pair $(x_i, m_i)$, answer for each query the
maximum value of $x_i \oplus \textit{nums}[j]$ over all $j$ with
$\textit{nums}[j] \le m_i$. If no element of $\textit{nums}$ satisfies
$\textit{nums}[j] \le m_i$, the answer for that query is $-1$.}

\begin{patternbox}
This is Section~\ref{sec:trie-max-xor}'s binary trie with one new
constraint layered on top: each query only considers a \textbf{subset}
of $\textit{nums}$ (those $\le m_i$), and the subset differs per query.
The keyword combination \emph{``maximum XOR'' plus ``subject to a
threshold that varies per query''} signals an \textbf{offline} technique:
sort the queries by their threshold and process them in increasing
order, incrementally growing the trie to include exactly the elements
eligible for each query, so that no element is ever inserted before it
is allowed to be, and no element needs to be removed once inserted.
\end{patternbox}

\subsection*{Understanding the problem carefully}

If queries arrived with arbitrary, unsorted thresholds $m_i$, we would
need to repeatedly figure out ``which elements of $\textit{nums}$ are
$\le m_i$'' for each query independently. But if we instead
\textbf{sort the queries by} $m_i$ and process them in increasing order,
the set of eligible elements only ever \textbf{grows} from one query to
the next (a higher threshold can only admit more elements, never fewer)
-- so we can maintain a single pointer into a separately sorted copy of
$\textit{nums}$, inserting new elements into the trie as the threshold
grows, and never needing to remove anything.

\begin{ideabox}[Offline: Sort Both Queries and the Array by Threshold]
Sort $\textit{nums}$ in increasing order. Sort the queries by $m_i$,
remembering each query's original index (to place answers back
correctly). Maintain a pointer $j$ into the sorted $\textit{nums}$,
starting at $0$. For each query, in increasing order of $m_i$: while $j$
has not exhausted $\textit{nums}$ and $\textit{nums}[j] \le m_i$, insert
$\textit{nums}[j]$ into the binary trie and advance $j$. Then, if the
trie is non-empty, answer this query using
Section~\ref{sec:trie-max-xor}'s $\textit{maxXorWith}(x_i)$; otherwise
(no eligible elements at all), answer $-1$.
\end{ideabox}

\subsection*{Why processing queries in sorted threshold order is correct and never wastes work}

Because eligibility only grows as the threshold increases, every element
inserted into the trie while answering an earlier (smaller-threshold)
query remains correctly eligible for every later (larger-threshold)
query too -- nothing inserted ever needs to be removed. This means the
pointer $j$ only ever moves forward, each element of $\textit{nums}$ is
inserted into the trie exactly once across the entire batch of queries,
and the total insertion work across all queries combined is $O(n)$
insertions, not $O(nq)$ for $q$ queries -- turning what looks like a
per-query rebuild into a single shared incremental structure.

\subsection*{Algorithm}

\begin{enumerate}
  \item Sort $\textit{nums}$ in increasing order.
  \item Sort the queries by $m_i$, keeping track of each query's
        original position.
  \item Initialize an empty binary trie and a pointer $j \gets 0$.
  \item For each query $(x_i, m_i)$, in sorted order:
  \begin{enumerate}
    \item While $j<n$ and $\textit{nums}[j] \le m_i$: insert
          $\textit{nums}[j]$ into the trie; $j \gets j+1$.
    \item If the trie is empty: this query's answer is $-1$.
    \item Else: this query's answer is
          $\textit{maxXorWith}(x_i)$.
  \end{enumerate}
  \item Place each computed answer back at its query's original
        position; return the answers.
\end{enumerate}

\subsection*{Dry run}

\dryrun{Take $\textit{nums} = [0,1,2,3,4]$ and queries
$[(3,1), (1,3), (5,6)]$.}

Sorted $\textit{nums}$ unchanged: $[0,1,2,3,4]$. Sorted queries by
$m_i$: $(3,1), (1,3), (5,6)$ (already in this order).

\begin{itemize}
  \item Query $(x{=}3, m{=}1)$: insert $\textit{nums}[j]$ while
        $\le1$: insert $0$ ($j\to1$), insert $1$ ($j\to2$); next is
        $2>1$, stop. Trie has $\{0,1\}$.
        $\textit{maxXorWith}(3)$: $3\oplus0=3$, $3\oplus1=2$; best is
        $3$.
  \item Query $(x{=}1, m{=}3)$: insert while $\le3$: insert $2$
        ($j\to3$), insert $3$ ($j\to4$); next is $4>3$, stop. Trie has
        $\{0,1,2,3\}$. $\textit{maxXorWith}(1)$: $1\oplus0=1$,
        $1\oplus1=0$, $1\oplus2=3$, $1\oplus3=2$; best is $3$.
  \item Query $(x{=}5, m{=}6)$: insert while $\le6$: insert $4$
        ($j\to5$); $j=n$, stop. Trie has $\{0,1,2,3,4\}$.
        $\textit{maxXorWith}(5)$: best pairing is $5\oplus2=7$ (checking
        all: $5{=}101$, $0{=}000\to101{=}5$, $1{=}001\to100{=}4$,
        $2{=}010\to111{=}7$, $3{=}011\to110{=}6$, $4{=}100\to001{=}1$);
        best is $7$.
\end{itemize}

Answers in original query order: $[3,3,7]$.

\subsection*{C++ implementation}

\begin{lstlisting}
#include <bits/stdc++.h>
using namespace std;

struct TrieNode {
    TrieNode* children[2] = {nullptr, nullptr};
};

class BinaryTrie {
    TrieNode* root = new TrieNode();
    static const int BITS = 32;
    bool empty = true;

public:
    void insert(int num) {
        empty = false;
        TrieNode* node = root;
        for (int b = BITS - 1; b >= 0; b--) {
            int bit = (num >> b) & 1;
            if (!node->children[bit]) node->children[bit] = new TrieNode();
            node = node->children[bit];
        }
    }

    bool isEmpty() { return empty; }

    int maxXorWith(int num) {
        TrieNode* node = root;
        int result = 0;
        for (int b = BITS - 1; b >= 0; b--) {
            int bit = (num >> b) & 1;
            int desired = 1 - bit;
            if (node->children[desired]) {
                result |= (1 << b);
                node = node->children[desired];
            } else {
                node = node->children[bit];
            }
        }
        return result;
    }
};

vector<int> maximizeXor(vector<int>& nums, vector<vector<int>>& queries) {
    sort(nums.begin(), nums.end());

    int q = queries.size();
    vector<int> order(q);
    for (int i = 0; i < q; i++) order[i] = i;
    sort(order.begin(), order.end(), [&](int a, int b) {
        return queries[a][1] < queries[b][1];
    });

    BinaryTrie trie;
    vector<int> answers(q);
    int j = 0, n = nums.size();

    for (int idx : order) {
        int x = queries[idx][0], m = queries[idx][1];
        while (j < n && nums[j] <= m) {
            trie.insert(nums[j]);
            j++;
        }
        answers[idx] = trie.isEmpty() ? -1 : trie.maxXorWith(x);
    }
    return answers;
}

int main() {
    vector<int> nums = {0, 1, 2, 3, 4};
    vector<vector<int>> queries = {{3,1}, {1,3}, {5,6}};

    for (int ans : maximizeXor(nums, queries)) cout << ans << " ";
    cout << endl;
    return 0;
}
\end{lstlisting}

\begin{complexitybox}
\textbf{Time Complexity.} Sorting $\textit{nums}$ and the queries costs
$O(n \log n + q \log q)$. Each element of $\textit{nums}$ is inserted
into the trie exactly once across the whole batch, costing
$O(n \cdot \textit{BITS})$ total. Each query performs one
$O(\textit{BITS})$ trie walk. Overall:
\[
  O((n+q)\log(n+q) + (n+q)\cdot \textit{BITS}).
\]
\textbf{Space Complexity.} $O(n \cdot \textit{BITS})$ for the trie, plus
$O(q)$ for the answers and sort order.
\end{complexitybox}

\begin{takeawaybox}
This closing problem of the chapter -- and of this entire multi-chapter
project -- is a deliberate capstone combining two ideas: the binary trie
from Section~\ref{sec:trie-max-xor}, and the classic \textbf{offline
query} technique of sorting queries by their constraint and processing
them alongside a sorted version of the data, so that a single
incrementally-built structure serves every query without ever needing to
be rebuilt or have elements removed. Whenever queries each restrict
attention to a growing prefix of sorted data, check whether sorting the
queries themselves by that same threshold turns many independent queries
into one shared incremental pass.
\end{takeawaybox}

